\chapter{总结与展望}

\section{研究工作总结}

回顾整个研究过程，本文最初想解决的问题很明确：C/C++库的模糊驱动太依赖人工编写，而现有的自动生成方法又受限于反馈粒度粗、单模型职责过载、缺乏知识积累和策略单一这几个瓶颈。为此本文设计并实现了MAGIC4FDG系统，尝试用多智能体协作加覆盖率反馈迭代的方式来突破这些限制。实验结果表明，两个核心问题："多智能体协作能否提升驱动质量"和"迭代反馈能否突破单轮瓶颈"，都得到了肯定的回答。下面是主要贡献的梳理：

（1）\textbf{提出了多智能体协作的模糊驱动生成架构}。本文把驱动生成拆成知识提取、策略规划、代码生成、编译修复、覆盖率评估、状态持久化和缺口诊断七个子任务，分别由专职智能体（Knowledge、Planner、Generation、Patching、Coverage、Checkpoint、Analyst）负责，通过LangGraph状态图做有向无环图调度。从数据来看，与相同LLM的单轮基线相比，这套架构把平均行覆盖率从23.97\%拉到了55.01\%。这组对比说明，职责分离确实能释放大语言模型在驱动生成上的潜力。

（2）\textbf{设计了"探索-利用"两阶段迭代机制}。第一轮通过LLM场景推理给每个槽位分配策略，快速建立覆盖率基线；第二轮起根据Analyst的结构化诊断做针对性改进。消融实验里，去掉迭代机制后覆盖率相对下降了25.2\%，这是所有组件中跌幅最大的，说明多轮反馈才是系统性能优势的根基。

（3）\textbf{实现了细粒度覆盖率反馈与知识累积}。Coverage智能体用LLVM工具链收集行级覆盖率，再结合CFG可达性分析过滤掉不可达代码；Analyst智能体把覆盖缺口诊断为结构化约束（未覆盖代码簇、触发条件、正负模式），并按槽位隔离累积。消融数据显示，去掉知识提取和缺口诊断分别导致18.5\%和13.8\%的相对下降，结构化知识对生成质量的支撑作用很明显。

（4）\textbf{构建了多策略并行生成框架}。系统内置10种策略（parse-centric、roundtrip、multi-api-sequence等），每轮多个槽位并行跑。实验中可以观察到不同策略在不同库上各有所长：cjson靠parse-centric拿到最高分，mbedtls则是differential表现最好。联合覆盖率远超任何单策略的最佳值，策略多样性的价值得到了验证。

（5）\textbf{在覆盖率维度取得了显著优势}。12个库的平均行覆盖率55.01\%，高于PromptFuzz\cite{lyu2024}的53.25\%和Hopper\cite{yu2024}的50.19\%。在相同的实验条件下，多智能体协作架构生成的驱动质量更高，覆盖率更优，说明精细化的职责分离和覆盖率反馈迭代确实能释放LLM在驱动生成上的潜力。

（6）\textbf{长时模糊测试验证了驱动的实际安全价值}。将生成的驱动在4个代表性库上执行3小时模糊测试，覆盖率均有显著提升（mbedtls从23.70\%提升至45.24\%），共触发9个崩溃并关联7个已知CVE（涵盖cjson、libpng、mbedtls三个库），还在cjson中发现了一个潜在新缺陷。这说明系统生成的驱动不仅在短时评估中表现优异，在实际安全测试场景中也具备发现真实漏洞的能力。

\section{局限性分析}

客观而言，MAGIC4FDG还有不少短板需要正视：

（1）\textbf{复杂状态依赖的驱动生成能力不足}。在harfbuzz、libxml2、libjpeg-turbo这些大型库上，系统的覆盖率和OSS-Fuzz官方的持续测试结果还有明显差距。本文认为最大的瓶颈在于：这些库的深层路径涉及复杂的状态机转换和多步协议交互，仅靠静态知识加覆盖率反馈，LLM很难推断出正确的调用序列。

（2）\textbf{LLM生成的不确定性}。同样的提示跑两次，产出的代码质量可能差很多。多策略并行在一定程度上对冲了这种随机性，但系统整体表现终究受限于LLM的能力天花板。

（3）\textbf{Token消耗与成本}。多智能体架构意味着多次LLM调用，单库迭代消耗约3M-24M tokens，取决于代码规模。资源有限时，覆盖率提升和计算成本之间需要做取舍。

（4）\textbf{环境依赖与泛化性}。系统目前绑定Docker加Clang/LLVM工具链。碰到非标准构建系统或特殊运行时环境的库，还需要额外适配。

（5）\textbf{语义正确性验证不足}。当前只通过编译成功和模糊测试运行来验证驱动的基本正确性，缺少对语义合理性的深层检查。有些驱动可能用不合理的API调用序列"骗"到了覆盖率，但并不代表有意义的测试场景。

\section{未来研究展望}

基于上面的局限性，本文认为有几个方向值得继续探索：

（1）\textbf{自适应策略选择与进化}（针对局限1、2）。目前10种策略是固定的，一个更有前景的方向是引入强化学习\cite{sutton2018}来做动态调整，根据目标库特征和历史反馈实时调整策略权重，让系统自己学会对不同库用不同打法。

（2）\textbf{长时间模糊测试与语料库管理}（针对局限1）。本文的长时测试实验已初步验证了延长时间预算的价值（mbedtls覆盖率从23.70\%提升至45.24\%），但目前只在4个库上做了验证，且时间预算固定为3小时。一个更有前景的方向是分层时间预算：覆盖率增长快的变体多给时间，同时做跨轮次的语料库继承和精简，让后续轮次能站在前面的肩膀上继续深挖。

（3）\textbf{多模型协作与成本优化}（针对局限3）。不是所有任务都需要最强的模型。编译修复这种相对机械的工作完全可以交给轻量级模型，把大模型的预算留给策略规划和代码生成，这样能在保持质量的同时压低Token开销。

（4）\textbf{跨项目知识迁移}（针对局限1、4）。不同库之间其实存在共性的API使用模式。如果能把在一个库上学到的有效模式（比如某类错误处理的写法、资源管理的套路）迁移到相似的库上，冷启动阶段的探索成本就能大幅降低。

（5）\textbf{与符号执行的深度融合}（针对局限1、5）。对于模糊测试难以穿透的复杂条件分支，符号执行技术\cite{cadar2008}可以辅助生成满足路径约束的输入构造代码。这和覆盖率引导的模糊测试形成互补，同时路径约束求解还能验证驱动的语义合理性，一举两得地突破深层路径和语义验证的双重瓶颈。

（6）\textbf{大规模持续集成部署}。把MAGIC4FDG接入CI/CD流水线，代码提交时自动为新增或修改的API生成驱动，这是系统走向工程实用的必经之路，也是后续最值得推进的方向。
